\section{Four hypotheses, stated so they can be falsified}
\label{sec:hypotheses}

\begin{description}[leftmargin=1.5em,style=nextline]
  \item[\textbf{H1 (Markovianity).}]
    The joint inference problem over a window --- or over a whole reanalysis ---
    is replaced by a \emph{sequence of simpler problems}, each conditionally
    independent of the past given the current state. This is what licenses
    cycling. Filtering takes observations up to $t$; 4D-Var and smoothers widen
    the unit to one window but keep the recursion between windows.
  \item[\textbf{H2 (model fidelity is the route to skill).}]
    The better the dynamical model, the better the assimilation.
  \item[\textbf{H3a (ODE/PDE state representation).}]
    The state is a point in a phase space whose prior is the flow map of a
    differential equation; uncertainty is carried as a perturbation of that
    point.
  \item[\textbf{H3b (autoregressive error propagation).}]
    Because that prior is generated by \emph{integrating} the model, every prior
    error is carried forward by the same recursion that carries the signal, and
    compounds along the window at the system's Lyapunov rate. Prior information
    reaches the analysis by one route only: through the integration.
\end{description}

\begin{table}[t]
\centering
\small
\begin{tabular}{@{}p{0.06\textwidth}p{0.26\textwidth}p{0.24\textwidth}p{0.32\textwidth}@{}}
\toprule
& \textbf{asserts} & \textbf{buys} & \textbf{costs} \\
\midrule
\textbf{H1} & the joint problem factorises into a sequence of conditionally
independent simpler ones & $O(1)$ memory, operational streaming, no training set
& cannot adapt to the \emph{observation configuration}; no amortisation across
windows \\[0.4em]
\textbf{H2} & skill is bounded by, and increasing in, model fidelity & the model
is a physically consistent, transferable prior & the model \emph{is} the prior,
so model error enters with no buffer \\[0.4em]
\textbf{H3a} & state $=$ point in ODE phase space; uncertainty $=$ perturbation
& tractable propagation; physical interpretability & the choice of state
variable fixes the conditioning of the inverse problem; a Gaussian update
implies $\PsiNG \equiv 0$ \\[0.4em]
\textbf{H3b} & prior errors propagate by the same integration as the signal & one
mechanism for signal and uncertainty alike & error compounds at the Lyapunov
rate, and \emph{model} sensitivity is inherited by the analysis whether or not
the \emph{posterior} has it \\
\bottomrule
\end{tabular}
\caption{The four hypotheses, what each buys, and what each costs.}
\label{tab:hypotheses}
\end{table}

\paragraph{H2 is load-bearing, and it is not simply false.}
The hypothesis is true as stated: better models do assimilate better. Our claim
is that it is \emph{binding}. Because the model supplies the entire prior, DA has
no mechanism to absorb model discrepancy, and the penalty is not a constant error
floor. No practitioner believes fidelity is the \emph{only} route to skill ---
model-error covariances, inflation and weak-constraint formulations all exist
precisely because model error is understood. The claim is that the available
buffers are weak, which Sections~\ref{sec:price} and \ref{sec:marginal} measure
rather than assert.

\paragraph{H3b is what makes H2 bind.}
H2 and H3b are not independent. Section~\ref{sec:sensitivity} shows that the
apparent force of ``better model $\Rightarrow$ better DA'' is partly an artefact
of the representation: it is the autoregressive propagation that routes all
prior information through the integration, and so makes the analysis inherit the
model's Lyapunov-amplified sensitivity whether or not the posterior has it.

\paragraph{A fifth hypothesis, named but not argued here.}
\textbf{H4 (gradients are valuable when available)} --- that the adjoint of the
dynamics, or of a variational cost built on it, is what makes an estimator
efficient. Its sharp form is a statement about the learned operator family and is
the subject of a companion paper: that \emph{differentiability is a training-time
requirement, not an inference-time one}. We retain here only its classical
corollary, which lives on the model-error axis
(Section~\ref{sec:adjoint}).
